% ---
% id: receiver/detector-imperfections/rem:effrank
% kind: rem
% title: "A port Gram whose rank drops under within-port attenuation"
% summary: >-
%   A port of two modes with Gram diag(1, ε²) has rank two; attenuating one of
%   its modes to zero leaves rank one. Apart from a zero efficiency,
%   within-port attenuation is the one operation in this contract that can
%   change the answer to the one-mode-per-port rank test.
% prominence: standard
% status: accepted
% uses: [receiver/thm:representation-independence, receiver/rem:rankone]
% mentions: [thm:efficiency, receiver/req:one-mode-per-port]
% assumes: []
% refs: []
% provenance: newly-derived
% lean:
% ---

\begin{remark}\label{rem:effrank}
Take $|\Inp|=2$ and a port collecting two modes, with the two rows of
$\Emd{\{k\}}\Gmat_{\Inp}$ equal to $(1,0)$ and $(0,\varepsilon)$ for some
$\varepsilon\in(0,1]$. Then $A_k=\operatorname{diag}(1,\varepsilon^2)$ has rank
$2$. Attenuating the second of those modes to zero, by a diagonal $D$ with
entries $1$ and $0$ inside the same port, gives
$A_k^{\mathrm{eff}}=\operatorname{diag}(1,0)$, of rank $1$. So a device that
does \emph{not} satisfy the one-mode-per-port criterion of
Theorem~\ref{thm:representation-independence}\,(iv) satisfies it after a within-port attenuation, and
the rank test gives its answer only on the Grams after every efficiency is
applied, never before (Requirement~\ref{req:one-mode-per-port}). For
single-mode ports the situation never arises, which is why the layer L1
declares one $\eta_k$ per port (the uniform case of
Theorem~\ref{thm:efficiency}\,(i)) by construction. Apart from declaring some
$\eta_k=0$, which sends $A_k$ to $0$ by Theorem~\ref{thm:efficiency}\,(i), this
is the one operation in this contract that can change the answer to the rank
test (Remark~\ref{rem:rankone}).
\end{remark}
